\documentclass[german,10pt]{book}      \input{../preambel}         \begin{document}  \setcounter{chapter}{0}

\chapter{Bell'sche Ungleichungen und Quantenkontextualit\"at}
% Kap x
\label{chap_Bell}

\info{Thomas Filk}{6.10.2024}%
Dies ist ein zweiter Kurztext zu Bell'schen Ungleichungen (vgl.\ \hyperref[chap_EPR]{EPR -- Bell -- CHSH}),
bei dem die Herleitung und die Schlussfolgerungen aus einer anderen Perspektive betrachtet
werden. In diese Betrachtung k\"onnen auch das Kochen-Specker Theorem und das
magische Quadrat von Peres und Mermin einbezogen werden.

Der allgemeine Formalismus besteht aus drei Schritten:
\begin{enumerate}
\item
Die Aufstellung einer Tabelle mit s\"amtlichen klassisch m\"oglichen Ergebnissen f\"ur
alle Observablen, die bei dem Problem betrachtet werden.
\item
Die \"Uberpr\"ufung einer Ungleichung oder Einschr\"ankung, die in allen erlaubten F\"allen
in der Tabelle g\"ultig ist. Unter der Annahme der Realit\"at muss diese Ungleichung bzw.\
Einschr\"ankung also f\"ur jedes System erf\"ullt sein.
\item
Der Nachweis, dass es Quantensysteme gibt, bei denen die Ungleichung bzw.\ die
Einschr\"ankungen verletzt sind. 
\end{enumerate}
Diese drei Schritte werden zun\"achst an einem Beispiel beschrieben. Anschlie\ss end erfolgt eine 
ausf\"uhrlichere allgemeine Beschreibung. Die sp\"ateren Abschnitte enthalten weitere
Beispiele. 

\section{Ein Beispiel: CHSH}

Die CHSH-Ungleichung bezieht sich auf ein System, das aus zwei Teilsystemen besteht.
Als Quantensystem betrachtet man meist zwei (in einem EPR-Zustand) verschr\"ankte Photonen, die sich
in gro\ss em Abstand befinden. An diesen beiden Photonen k\"onnen die Polarisationen unter
verschiedenen Winkeln gemessen werden. 

Bei der CHSH-Ungleichung, die auch von Alain Aspect experimentell getestet
wurde, handelt es sich um vier Observablen $\{A_1, A_2, B_1, B_2\}$, von denen $A_1$ und $A_2$ nicht
gleichzeitig gemessen werden k\"onnen (sie geh\"oren zu verschiedenen Polarisationsrichtungen an einem
Photon, die relativ zueinander um $45^\circ$ gedreht sind) und ebenso sind $B_1$ und $B_2$ nicht 
kompatibel. Das bedeutet, nur die vier paarweisen Kombinationen $(A_i,B_j)$ ($i,j=1,2$) k\"onnen gleichzeitig, 
gemessen werden. Au\ss erdem k\"onnen die Ergebnisse f\"ur jede Observable nur $+1$ oder $-1$ sein. 

\subsection{Die Realit\"atsannahme  -- Die klassische Tabelle}

F\"ur vier Observable, die nur jeweils zwei Werte ($+1$ und $-1$) annehmen k\"onnen, gibt es insgesamt
$2^4=16$ M\"oglichkeiten, die in Tabelle \ref{tab_CHSH} zusammengefasst werden. Unter der Annahme
der Realit\"at muss f\"ur jedes einzelne System eine dieser 16 M\"oglichkeiten realisiert sein. 

\begin{SCtable}[50][htb]
\begin{tabular}{cccc|c||}
$A_1$ & $A_2$ & $B_1$ & $B_2$ & $S$ \\ \hline
$+1$ & $+1$ & $+1$ & $+1$ &  $+2$ \\
$+1$ & $+1$ & $+1$ & $-1$ &  $+2$ \\
$+1$ & $+1$ & $-1$ & $+1$ &  $-2$ \\
$+1$ & $+1$ & $-1$ & $-1$ &  $-2$ \\
$+1$ & $-1$ & $+1$ & $+1$ &  $+2$ \\
$+1$ & $-1$ & $+1$ & $-1$ &  $-2$ \\
$+1$ & $-1$ & $-1$ & $+1$ &  $+2$ \\
$+1$ & $-1$ & $-1$ & $-1$ &  $-2$ \\
\end{tabular}
\begin{tabular}{cccc|c}
$A_1$ & $A_2$ & $B_1$ & $B_2$ & $S$ \\ \hline
$-1$ & $+1$ & $+1$ & $+1$ &  $-2$ \\
$-1$ & $+1$ & $+1$ & $-1$ &  $+2$ \\
$-1$ & $+1$ & $-1$ & $+1$ &  $-2$ \\
$-1$ & $+1$ & $-1$ & $-1$ &  $+2$ \\
$-1$ & $-1$ & $+1$ & $+1$ &  $-2$ \\
$-1$ & $-1$ & $+1$ & $-1$ &  $-2$ \\
$-1$ & $-1$ & $-1$ & $+1$ &  $+2$ \\
$-1$ & $-1$ & $-1$ & $-1$ &  $+2$ 
\end{tabular}
\caption{\label{tab_CHSH}%
F\"ur die vier Observable $A_1,A_2,B_1,B_2$ sind nur die angegebenen 16 Kombinationen
m\"oglich. In allen 16 F\"allen hat die Gr\"o\ss e $S=a_1b_1+a_1b_2+a_2b_1-a_2b_2$
den Wert $+2$ oder $-2$. F\"ur ein beliebiges Ensemble aus den angegebenen M\"oglichkeiten
muss der Erwartungswert von $S$ daher immer zwischen $+2$ und $-2$ liegen bzw.\
$|\langle S \rangle | \leq 2$. }
\end{SCtable}

\subsection{Die Ungleichung}





\section{Der allgemeine Formalismus}


\subsection{Die Annahme der Realit\"at}

Ausgangspunkt ist eine Menge von Observablen $\{A,B,C,...\}$, die nicht alle gleichzeitig
gemessen werden k\"onnen, d.h., die nicht alle miteinander kommutieren. Es gibt jedoch bestimmte
Teilmengen von Observablen, die miteinander kompatibel sind und somit gleichzeitig gemessen
werden k\"onnen. Jede der Observablen kann nur bestimmte Werte annehmen -- in den meisten 
F\"allen handelt es sich um bin\"are Werte, d.h., eine Messung der Observablen kann nur einen
von zwei m\"oglichen Messwerten ergeben. 

Bei der CHSH-Ungleichung, die auch von Alain Aspect experimentell getestet
wurde, handelt es sich um vier Observable $\{A_1, A_2, B_1, B_2\}$, von denen $A_1$ und $A_2$ nicht
gleichzeitig gemessen werden k\"onnen (sie geh\"oren zu verschiedenen Polarisationsrichtungen,
die relativ zueinander um $45^\circ$ gedreht sind) und ebenso sind $B_1$ und $B_2$ nicht kompatibel.
Das bedeutet, nur die paarweisen Kombinationen $(A_i,B_j)$ ($i,j=1,2$) k\"onnen gleichzeitig, gemessen
werden. Au\ss erdem k\"onnen die Ergebnisse f\"ur jede Observable nur $+1$ oder $-1$ sein. 

Man kann nun eine Tabelle erstellen, in der alle m\"oglichen Messwerte von allen Observablen
aufgelistet sind. Diese Tabelle entspricht einer \glqq Realit\"atsannahme\grqq: Es wird angenommen,
dass jeder Observablen unabh\"angig davon, ob sie gemessen wird oder nicht, ein Wert zukommt,
der bei einer Messung auch ihrem Messwert entspricht. Unter dieser Annahme muss ein gegebenes
System einer der M\"oglichkeiten in der Tabelle entsprechen. Eventuell unterliegen die Observablen
gewissen einschr\"ankenden Bedingungen: Bei einem EPR-Zustand muss die Messung einer Observablen
an einem der Teilsysteme immer das entgegengesetzte Ergebnis von der Messung der entsprechenden
Observablen am anderen Teilsystem ergeben. Das bedeutet, die M\"oglichen Messwerte zu den
entsprechenden Observablen an den Teilsystemen m\"ussen immer entgegengesetzt sein. 
L\"asst man insgesamt sechs Observable zu -- $\{A_1,B_1,C_1,A_2,B_2,C_2\}$, wobei sich $A_1,B_1,C_1$ auf
die Messung der Polarisationsrichtung an Photon 1 unter bestimmten Winkeln beziehen und entsprechend
$A_2,B_2,C_2$ auf die Messungen derselben Polarisationsrichtungen an Photon 2 -- dann sind die
Werte in der Tabelle bez\"uglich gleicher Messungen immer antikorreliert. Von den zun\"achst $2^6=64$
M\"oglichkeiten, bei sechs Observablen jeweils die Messwerte $+1$ oder $-1$ zu finden, bleiben nur
$2^3=8$ M\"oglichkeiten \"ubrig, da die Messwerte von $A_2.B_2,C_2$ bereits durch die Messwerte von
$A_1,B_1,C_1$ festgelegt sind. 

\subsection{Die Ungleichung}

Die schwierige Aufgabe besteht darin, eine Bedingung an bestimmte Korrelationen
in der Tabelle zu finden, die von allen erlaubten M\"oglichkeiten erf\"ullt werden, zu denen es
aber in der Quantentheorie Systeme gibt, bei denen diese Bedingungen verletzt sind. 

In vielen F\"allen handelt es sich um eine Ungleichung, die von einer bestimmten
Kombination von Korrelationen erf\"ullt ist. In manchen F\"allen, beispielsweise beim magischen
Quadrat von Peres und Mermin, l\"asst sich eine bestimmte Bedingung klassisch \"uberhaupt nicht
erf\"ullen, wohingegen die Korrelationen an verschr\"ankten Quantensystemen die
Bedingungen erf\"ullen. Es gibt in vielen F\"allen keine Regel, nach der eine solche
Bedingung gefunden werden kann, doch in den meisten F\"allen ist diese Regel bekannt
und man muss sich nur anhand der Tabelle \"uberzeugen, dass die Regel in allen F\"allen
erf\"ullt ist.

\subsection{Die Verletzung durch ein Quantensystem}











\end{document}

